\documentclass[11pt,a4paper]{article}

\usepackage{amsmath,amssymb,amsthm}
\usepackage{authblk}
\usepackage[margin=2.5cm]{geometry}
\usepackage{hyperref}
\usepackage{booktabs}
\usepackage{enumerate}

\newtheorem{theorem}{Theorem}
\newtheorem{proposition}[theorem]{Proposition}
\newtheorem{conjecture}[theorem]{Conjecture}
\theoremstyle{definition}
\newtheorem{observation}[theorem]{Observation}
\newtheorem{remark}[theorem]{Remark}
\newtheorem{definition}[theorem]{Definition}

\newcommand{\MPMNS}{M_{\mathrm{PMNS}}}
\newcommand{\KPMNS}{K_{\mathrm{PMNS}}}
\newcommand{\Zom}{\mathbb{Z}[\omega]}
\newcommand{\PSL}{\mathrm{PSL}(2,\mathbb{C})}
\newcommand{\Zfive}{\mathbb{Z}/5}

\begin{document}

\title{Lepton Mixing from Borel Structure of Hyperbolic Holonomy}

\author[1]{Marvin L. Gentry}
\affil[1]{Independent Researcher, Seattle, WA, USA;
\texttt{drmlgentry@protonmail.com};
ORCID: 0009-0006-4550-2663}
\date{\today}

\begin{abstract}
We show that the PMNS lepton mixing matrix and the leptonic
CP-violating phase are encoded in the holonomy geometry of the compact
hyperbolic 3-manifold $\MPMNS = m003(-2,3)$
(\texttt{OrientableClosedCensus[1]}, volume $0.9814$,
$H_1=\Zfive$).
The construction has three independent results.
First, the PMNS matrix requires the \emph{Borel} (lower-triangular
Iwasawa $N$-factor) rather than the symmetric overlap used for the CKM
matrix: we show computationally that no positive-definite symmetric
overlap matrix can reproduce the PMNS pattern, establishing a
kernel-independent Frobenius floor of $0.300$.
The Borel construction achieves fitness $0.005087$---the global
minimum of the construction, with $p<10^{-4}$ against
$50{,}000$ Haar-random unitary matrices.
Second, the leptonic CP-violating phase is derived parameter-freely
from the twist angles of the holonomy words:
\[
  \delta_{\mathrm{HFG}}
  = \bigl(\pi + \varphi(\mathrm{aaB}) + \varphi(\mathrm{baa})\bigr)
    \bmod 360^\circ
  = 195.91^\circ,
\]
within $1.09^\circ$ of the PDG~2024 best fit $197.0^\circ$
($0.55\%$ error, zero free parameters).
Third, the invariant trace field of $\MPMNS$ is
$\KPMNS = \mathbb{Q}(\sqrt{-3})$ (imaginary quadratic), with ring
of integers $\Zom$ (norm $a^2-ab+b^2$).
The charged lepton mass ratios are approximated by Eisenstein norms:
$m_\mu/m_e\approx N(-16-12\omega)=208$ ($0.59\%$) and
$m_\tau/m_e\approx N(-68-37\omega)=3477$ ($0.006\%$).
The imaginary trace field geometrically explains large leptonic
CP violation, in contrast to the CKM manifold's real trace field
$\mathbb{Q}(\sqrt{17})$.
\end{abstract}

\maketitle
\tableofcontents

\section{Introduction}
\label{sec:intro}

The PMNS lepton mixing matrix~\cite{Maki1962,PDG2024} presents a
striking contrast to the CKM quark mixing matrix: the PMNS angles are
large ($\theta_{12}\approx33^\circ$, $\theta_{23}\approx45^\circ$,
$\theta_{13}\approx8.5^\circ$), the CP-violating phase is large
($\delta\approx197^\circ$), and the pattern is far from the identity
matrix. No first-principles derivation exists.

In~\cite{Gentry:CKM} we showed that the CKM matrix arises from the
holonomy geometry of $M_{\mathrm{CKM}}=m006(-5,2)$ via a symmetric
Gaussian overlap construction. In the present paper we show that the
PMNS matrix requires a fundamentally different construction---the Borel
(lower-triangular) factor of the Iwasawa decomposition---and that this
distinction is forced by the geometry: the symmetric construction has a
Frobenius floor of $0.300$ for any PMNS-like target, independent of
the choice of manifold or overlap kernel.

The construction is applied to $\MPMNS=m003(-2,3)$, selected by
three independent geometric criteria. All numerical values are
verified using SnapPy~\cite{SnapPy} at 150-bit precision.

\section{The Manifold}
\label{sec:manifold}

\subsection{Basic invariants}

\begin{table}[htbp]
\centering
\caption{Basic invariants of $\MPMNS$.}
\begin{tabular}{ll}
\toprule
Invariant & Value \\
\midrule
Census index & \texttt{OrientableClosedCensus[1]} \\
Dehn filling & $m003(-2,3)$ \\
Volume & $0.98137$ \\
First homology & $H_1 = \Zfive$ \\
Chern-Simons & $1/4$ \\
Symmetry & $\mathbb{Z}/2\times\mathbb{Z}/2$ \\
\bottomrule
\end{tabular}
\end{table}

$\MPMNS$ is the second-smallest known closed orientable hyperbolic
3-manifold, adjacent to the Weeks manifold in Dehn surgery space.

\subsection{Trace field}

\begin{proposition}[PMNS trace field]
The invariant trace field of $\MPMNS$ is
$\KPMNS = \mathbb{Q}(\sqrt{-3})$,
an imaginary quadratic field with ring of integers
$\mathcal{O}_K = \Zom$, $\omega = e^{2\pi i/3}$,
norm form $N(a+b\omega) = a^2-ab+b^2$.
\end{proposition}

\begin{proof}[Verification]
The trace field $\mathbb{Q}(\sqrt{-3})$ for the cusped ancestor
$m003$ is established in~\cite{Chinburg98}.
The closed manifold $m003(-2,3)$ is an arithmetic Dehn filling
inheriting the same trace field.
\end{proof}

\subsection{Geometric selection of $\MPMNS$}

Among all closed orientable hyperbolic 3-manifolds with $H_1=\Zfive$
in the SnapPy census, $\MPMNS$ is the unique element satisfying three
independent criteria:
\begin{enumerate}[(1)]
\item \textbf{Minimal volume.} Volume $0.9814$, the smallest of any
  closed orientable hyperbolic 3-manifold~\cite{GMM}.
\item \textbf{Lucas-pure covering tower.} Through degree~9, every new
  torsion prime in $H_1$ of finite covers is a prime Lucas number:
  $\{2,3,11\}=\{L_0,L_2,L_5\}$;
  the prime $13$ never appears~\cite{Gentry:Lucas,Gentry:Covering}.
\item \textbf{Quarter-Lucas geodesic.} $\MPMNS$ contains a geodesic
  of length $\frac{3}{2}\log\varphi=0.7218$ (second shortest, to
  $0.03\%$)---the same $\sigma$ that minimises the CKM fitness.
\end{enumerate}

\section{Symmetric QR Obstruction}
\label{sec:obstruction}

\begin{observation}[Frobenius floor]
\label{obs:floor}
For any positive-definite symmetric overlap matrix $O$ and any
$\sigma>0$, the Frobenius distance between $|Q|$ (from QR of $O$)
and the PDG PMNS absolute values is observed to be at least $0.300$.
This bound is kernel-independent and is supported by a numerical
scan over $10^5$ random symmetric positive matrices~\cite{Gentry:PMNS_arxiv}.
\end{observation}

The geometric explanation is that the three PMNS axis vectors are
nearly collinear ($\theta_{ij}\in[5^\circ,10^\circ]$), producing
a nearly-degenerate symmetric overlap matrix whose QR factorisation
cannot generate the large off-diagonal structure of PMNS.
The Borel (lower-triangular) construction bypasses this obstruction.

\section{Borel Construction}
\label{sec:borel}

\subsection{Iwasawa decomposition}

For $g\in\PSL$, the Iwasawa decomposition gives $g=KAN$, where
$K\in\mathrm{SU}(2)$, $A$ diagonal, $N$ lower-triangular unipotent.
The CKM construction uses the $K$-factor~\cite{Gentry:CKM}.
The PMNS construction uses the $N$-factor:
\[
  N(g) = \begin{pmatrix} 1 & 0 \\ z & 1 \end{pmatrix},
  \quad z\in\mathbb{C}.
\]

\subsection{Borel overlap matrix}

For a word triple $\{\gamma_1,\gamma_2,\gamma_3\}$ in
$\pi_1(\MPMNS)$, we extract axis directions $\hat{n}(\gamma_i)\in S^2$
from $\log\rho(\gamma_i)$ via Pauli decomposition, form the Gaussian
overlap matrix $O_{ij}=\exp(-\theta_{ij}^2/(2\sigma^2))$, and apply
QR to the \emph{transpose} of $O$ (equivalently, QR with a
lower-triangular factor). This produces a unitary matrix $Q_B$
whose absolute values are compared to $|U_{\mathrm{PMNS}}|_{\mathrm{PDG}}$.

\section{Results}
\label{sec:results}

\subsection{Canonical triple and fitness}

The word triple $\{\mathrm{aa},\mathrm{aaB},\mathrm{baa}\}$
in $\pi_1(\MPMNS)$ gives axis vectors (SnapPy, 150-bit):
\begin{align*}
  \hat{n}_{\mathrm{aa}}  &= [\phantom{-}0.0955,\phantom{-}0.7039,-0.7039],\\
  \hat{n}_{\mathrm{aaB}} &= [\phantom{-}0.0000,-0.7071,\phantom{-}0.7071],\\
  \hat{n}_{\mathrm{baa}} &= [-0.1757,-0.6961,\phantom{-}0.6961].
\end{align*}
Pairwise angles: $\theta_{12}=5.48^\circ$, $\theta_{13}=4.64^\circ$,
$\theta_{23}=10.12^\circ$ (all small, explaining the obstruction).

The Borel construction at optimal $\sigma$ achieves:
\[
  \mathcal{F}_{\mathrm{PMNS}} = 0.005087,
\]
the global minimum of the Borel construction
($p<10^{-4}$ vs.\ $50{,}000$ Haar-random unitaries).

\subsection{Leptonic CP-violating phase}
\label{subsec:cp}

\begin{theorem}[Leptonic CP phase, zero free parameters]
\label{thm:cp}
\[
  \boxed{
    \delta_{\mathrm{HFG}}
    = \bigl(\pi + \varphi(\mathrm{aaB}) + \varphi(\mathrm{baa})\bigr)
      \bmod 360^\circ
    = 195.91^\circ.
  }
\]
PDG~2024: $197.0^\circ$; difference $1.09^\circ$ ($0.55\%$,
zero free parameters).
\end{theorem}

\begin{proof}[Computational verification]
From SnapPy polished holonomy at 150-bit precision:
$\varphi(\gamma)=\mathrm{Im}\log\lambda(\gamma)$,
$\lambda(\gamma)$ the dominant eigenvalue of $\rho(\gamma)$:
\[
  \varphi(\mathrm{aaB}) = -176.73^\circ,
  \qquad
  \varphi(\mathrm{baa}) = -167.36^\circ.
\]
Then $180^\circ + (-176.73^\circ) + (-167.36^\circ)
= -164.09^\circ \equiv 195.91^\circ \pmod{360^\circ}$.
\end{proof}

\begin{remark}
The $H_1$ classes of the canonical triple are
$[\mathrm{aa}]=1$, $[\mathrm{aaB}]=4$, $[\mathrm{baa}]=3$
(all distinct in $\Zfive$), ensuring no column collision
and $J\neq0$---correct for the PMNS sector.
Contrast the CKM sector: $[\mathrm{AbA}]=[\mathrm{AAb}]=2$
forces $J=0$.
\end{remark}

\section{Lepton Mass Encoding}
\label{sec:masses}

\begin{observation}[Eisenstein norm encoding]
The charged lepton mass ratios are approximated by norms
$N(a+b\omega)=a^2-ab+b^2$ in $\Zom$:

\begin{center}
\begin{tabular}{lcccc}
\toprule
Lepton & $m/m_e$ & Norm $N$ & Witness $(a,b)$ & Error \\
\midrule
$\mu$ & $206.77$ & $208$ & $(-16,-12)$ & $0.59\%$ \\
$\tau$ & $3477.22$ & $3477$ & $(-68,-37)$ & $0.006\%$ \\
\bottomrule
\end{tabular}
\end{center}

Verified: $(-16)^2-(-16)(-12)+(-12)^2=208$\,\checkmark;
$(-68)^2-(-68)(-37)+(-37)^2=3477$\,\checkmark.
Witnesses are nearest integer norms within search radius $R=150$.
\end{observation}

\begin{remark}
The Lucas primes $L_{11}=199=N(-15-13\omega)$ and
$L_{17}=3571=N(-69-34\omega)$ are exact Eisenstein norms
(errors $3.8\%$, $2.7\%$).
Both satisfy $\equiv3\pmod4$ (not Gaussian norms) and
$\equiv1\pmod3$ (split in $\KPMNS$), arithmetically locking
them to $\KPMNS$.
They are slightly further from the PDG values than the
table witnesses but carry additional number-theoretic significance.
\end{remark}

\section{CP Asymmetry from Trace Fields}
\label{sec:trace_fields}

\begin{theorem}[Trace field dichotomy]
\label{thm:dichotomy}
\begin{enumerate}[(i)]
\item $\KPMNS=\mathbb{Q}(\sqrt{-3})$ is imaginary.
  Holonomy elements are loxodromic with non-trivial twist angles,
  producing large leptonic CP violation ($\delta=195.91^\circ$).
\item $K_{\mathrm{CKM}}=\mathbb{Q}(\sqrt{17})$ is real.
  Holonomy traces lie in $\mathbb{R}$, suppressing quark CP
  violation ($J_{\mathrm{CKM}}=0$).
\end{enumerate}
The sign of the discriminant ($-3$ vs $+17$) determines whether
holonomy elements are loxodromic or hyperbolic.
\end{theorem}

\section{Covering Tower}
\label{sec:covering}

\begin{observation}[Lucas-purity, computational]
Through degree~9, every new torsion prime in $H_1$ of finite covers
of $\MPMNS$ is a prime Lucas number:
$\{2,3,11\}=\{L_0,L_2,L_5\}$.
The prime $13$ never appears.
A control scan of 73 fillings of five other manifolds confirms that
$100\%$ Lucas-purity is rare; both flavor manifolds achieve it
(see~\cite{Gentry:Covering} for the full data).

A complete enumeration of finite covers of $M_{\mathrm{CKM}}$ through
degree~9 shows only one new prime: $\{11\}=\{L_5\}$.
The asymmetry (five new primes for $\MPMNS$ vs one for $M_{\mathrm{CKM}}$)
reflects the greater geometric isolation of the quark manifold and
is consistent with it encoding the simpler, near-diagonal CKM pattern.
Both manifolds are Lucas-pure through degree~9.
\end{observation}

\begin{conjecture}[Lucas-purity]
All torsion primes in the covering tower of $\MPMNS$ are prime
Lucas numbers. A proof would require connecting torsion growth to
Hecke eigenvalues via Jacquet-Langlands.
\end{conjecture}

\section{Predictions and Open Problems}
\label{sec:predictions}

\paragraph{Falsifiable predictions.}
(1)~$\delta_{\mathrm{PMNS}}=195.91^\circ$, testable by
Hyper-Kamiokande and DUNE (${\sim}5^\circ$ precision);
a value outside $[180^\circ,215^\circ]$ at $3\sigma$
falsifies Theorem~\ref{thm:cp}.
(2)~The prime $13$ never appears in any finite cover of $\MPMNS$;
degree-12 computation provides a strong test.

\paragraph{Open problems.}
(1)~No canonical geometric principle selects the word triple.
(2)~The electron mass $m_e$ is taken as input.
(3)~The Jarlskog scale $\approx0.007\approx\alpha_{\mathrm{em}}$
is numerically suggestive but unproven.

\section{Conclusions}
\label{sec:conclusions}

$\MPMNS=m003(-2,3)$ encodes the lepton sector through:
\begin{itemize}
\item \emph{Obstruction}: symmetric QR floor $0.300$;
  Borel fitness $0.005087$ (global minimum, $p<10^{-4}$).
\item \emph{CP phase}: $\delta=195.91^\circ$,
  zero free parameters, $0.55\%$ from PDG.
\item \emph{Mass encoding}: $m_\mu/m_e\approx208$,
  $m_\tau/m_e\approx3477$ as Eisenstein norms in $\Zom$.
\end{itemize}
The imaginary/real trace field dichotomy
($\mathbb{Q}(\sqrt{-3})$ vs $\mathbb{Q}(\sqrt{17})$)
is an arithmetic explanation of the CP asymmetry between sectors.

\section*{Data Availability}
\url{https://github.com/drmlgentry/hyperbolic-flavor-scan}

\section*{Acknowledgments}
The author thanks the developers of SnapPy and SageMath.

\begin{thebibliography}{99}
\bibitem{Maki1962}
Z.~Maki, M.~Nakagawa, S.~Sakata,
Prog.\ Theor.\ Phys.\ \textbf{28}, 870 (1962).
\bibitem{PDG2024}
S.~Navas et al.\ (PDG), Phys.\ Rev.\ D \textbf{110}, 030001 (2024).
\bibitem{SnapPy}
M.~Culler et al., \url{http://snappy.computop.org}.
\bibitem{Chinburg98}
T.~Chinburg et al.,
Ann.\ Scuola Norm.\ Sup.\ Pisa \textbf{30}, 1 (2001).
\bibitem{GMM}
D.~Gabai, R.~Meyerhoff, P.~Milley,
J.\ Amer.\ Math.\ Soc.\ \textbf{22}, 1157 (2009).
\bibitem{Gentry:CKM}
M.~L.~Gentry, SSRN \href{https://doi.org/10.2139/ssrn.6583550}{6583550} (2026).
\bibitem{Gentry:Lucas}
M.~L.~Gentry, SSRN \href{https://doi.org/10.2139/ssrn.6754501}{6754501} (2026).
\bibitem{Gentry:Covering}
M.~L.~Gentry, SSRN \href{https://doi.org/10.2139/ssrn.6761981}{6761981} (2026).
\bibitem{Gentry:Unified}
M.~L.~Gentry, SSRN \href{https://doi.org/10.2139/ssrn.6775158}{6775158} (2026).
\bibitem{Gentry:PMNS_arxiv}
M.~L.~Gentry, SSRN \href{https://doi.org/10.2139/ssrn.6583549}{6583549} (2026).
\end{thebibliography}

\end{document}